\documentclass[oneside]{book}
\usepackage[backend=biber,natbib=true,style=authoryear]{biblatex}
\addbibresource{/home/nguyen/1_NQBH/reference/bib.bib}
\usepackage[utf8]{inputenc}
\usepackage{float}
\usepackage{graphicx}
\usepackage[colorlinks=true,linkcolor=blue,urlcolor=red,citecolor=magenta]{hyperref}
\usepackage{amsmath,amssymb,amsthm}
\allowdisplaybreaks
\numberwithin{equation}{section}
\newtheorem{assumption}{Assumption}[section]
\newtheorem{lemma}{Lemma}[section]
\newtheorem{corollary}{Corollary}[section]
\newtheorem{definition}{Definition}[section]
\newtheorem{proposition}{Proposition}[section]
\newtheorem{theorem}{Theorem}[section]
\newtheorem{notation}{Notation}[section]
\newtheorem{remark}{Remark}[section]
\newtheorem{example}{Example}[section]
\newtheorem{ques}{Question}[section]
\newtheorem{problem}{Problem}[section]
\newtheorem{conjecture}{Conjecture}[section]
\usepackage[left=0.5in,right=0.5in,top=1.5cm,bottom=1.5cm]{geometry}
\usepackage{fancyhdr}
\pagestyle{fancy}
\fancyhf{}
\addtolength{\headheight}{0pt}% obsolete
\lhead{\scriptsize \chaptername~\thechapter}
\rhead{\scriptsize \nouppercase{\leftmark}} %\nouppercase !
\renewcommand{\chaptermark}[1]{\markboth{#1}{}}
\cfoot{\thepage}
\def\labelitemii{$\circ$}

\title{Health}
\author{Collector: Nguyen Quan Ba Hong}
\date{\today}

\begin{document}
\maketitle
\setcounter{secnumdepth}{6}
\setcounter{tocdepth}{6}
\tableofcontents

%------------------------------------------------------------------------------%

\chapter{Wikipedia's}

\section{\href{https://en.wikipedia.org/wiki/Crunch_(exercise)}{Wikipedia\texttt{/}Crunch (Exercise)}}
\textsf{Performing the crunch.}

The \textit{crunch} is 1 of the most popular \href{https://en.wikipedia.org/wiki/Abdominal_exercise}{abdominal exercises}.

When performed properly, it engages all the abdominal muscles but primarily it works the \href{https://en.wikipedia.org/wiki/Rectus_abdominis_muscle}{rectus abdominis muscle} and the \href{https://en.wikipedia.org/wiki/Abdominal_external_oblique_muscle}{obliques}.

It allows both building 6-pack abs, and tightening the belly.

Crunches use the exerciser's own body weight to tone muscle, and are recommended as a low-cost exercise that can be performed at home.[1]

\subsection{Form}
\textsf{In a crunch, the lower back does not lift off the floor.}

The \href{https://en.wikipedia.org/wiki/Biomechanics}{biomechanics} professor Stuart McGill was quoted in \textit{\href{https://en.wikipedia.org/wiki/The_New_York_Times}{The New York Times} Health} blog as stating:
\begin{quotation}
    An approved crunch begins with you lying down, 1 knee bent, and hands positioned beneath your lower back for support.
    
    ``Do not hollow your stomach or press your back against the floor'', McGill says.
    
    Gently lift your head and shoulders, hold briefly and relax back down.[2]
\end{quotation}
In a crunch, unlike a \href{https://en.wikipedia.org/wiki/Sit-up}{sit-up}, the \href{https://en.wikipedia.org/wiki/Lower_back}{lower back} stays on the floor.

This is said to eliminate any involvement by the \href{https://en.wikipedia.org/wiki/Hip_flexors}{hip flexors}, and make the crunch an effective \href{https://en.wikipedia.org/wiki/Isolation_exercise}{isolation exercise} for the abdominals.[3]

\subsection{See also}
\begin{itemize}
    \item \href{https://en.wikipedia.org/wiki/Plank_(exercise)}{Plank (exercise)}
    \item \href{https://en.wikipedia.org/wiki/Sit-up}{Sit-up}\hfill$\square$
\end{itemize}

%------------------------------------------------------------------------------%

\section{\href{https://en.wikipedia.org/wiki/Plank_(exercise)}{Wikipedia\texttt{/}Plank (Exercise)}}
\textsf{Recruit performing a plank at a US Coast Guard training.}

The \textit{plank} (also called a \textit{front hold, hover}, or \textit{abdominal bridge}) is an \href{https://en.wikipedia.org/wiki/Isometric_exercise}{isometric} \href{https://en.wikipedia.org/wiki/Core_(anatomy)}{core} strength exercise that involves maintaining a position similar to a \href{https://en.wikipedia.org/wiki/Push-up}{push-up} for the maximum possible time.

\subsection{Form}
The most common plank is the forearm plank which is held in a \href{https://en.wikipedia.org/wiki/Push-up}{push-up}-like position, with the body's weight borne on forearms, elbows, and toes.

Many variations exist such as the side plank and the reverse plank.[1][2]

The plank is commonly practiced in \href{https://en.wikipedia.org/wiki/Pilates}{Pilates} and \href{https://en.wikipedia.org/wiki/Yoga}{yoga}, and by those training for \href{https://en.wikipedia.org/wiki/Boxing}{boxing} and other sports.[3][4][5]

%
The ``extended plank'' adds substantial difficulty to the standard plank exercise.

To perform the extended plank, a person begins in the push-up position and then extends the arms or hands as far forward as possible.[6]

\subsection{Effect}
The plank strengthens the abdominals, back and shoulders.

Muscles involved in the front plank include:[7]
\begin{itemize}
    \item Primary muscles: \href{https://en.wikipedia.org/wiki/Erector_spinae_muscles}{erector spinae}, \href{https://en.wikipedia.org/wiki/Rectus_abdominis_muscle}{rectus abdominis} (abs), and \href{https://en.wikipedia.org/wiki/Transversus_abdominis_muscle}{transverse abdominis}.
    \item Secondary muscles (\href{https://en.wikipedia.org/wiki/Anatomical_terms_of_muscle#Synergist}{synergists}/\href{https://en.wikipedia.org/wiki/Segmental_stabilizers}{segmental stabilizers}): \href{https://en.wikipedia.org/wiki/Trapezius_muscle}{trapezius} (traps), \href{https://en.wikipedia.org/wiki/Rhomboid_muscles}{rhomboids}, \href{https://en.wikipedia.org/wiki/Rotator_cuff}{rotator cuff}, the anterior, medial, and posterior \href{https://en.wikipedia.org/wiki/Deltoid_muscle}{deltoid muscles} (delts), \href{https://en.wikipedia.org/wiki/Pectoralis_major_muscle}{pectorals} (pecs), \href{https://en.wikipedia.org/wiki/Serratus_anterior_muscle}{serratus anterior}, \href{https://en.wikipedia.org/wiki/Gluteus_maximus_muscle}{gluteus maximus} (glutes), \href{https://en.wikipedia.org/wiki/Quadriceps_femoris_muscle}{quadriceps} (quads), and \href{https://en.wikipedia.org/wiki/Gastrocnemius_muscle}{gastrocnemius}.
\end{itemize}
Muscles involved in the side plank include: 
\begin{itemize}
    \item Primary: \href{https://en.wikipedia.org/wiki/Transversus_abdominis_muscle}{transversus abdominis muscle}, \href{https://en.wikipedia.org/wiki/Gluteus_medius_muscle}{gluteus medius} and \href{https://en.wikipedia.org/wiki/Gluteus_minimus_muscle}{gluteus minimus} muscles (abductors), the \href{https://en.wikipedia.org/wiki/Adductor_muscles_of_the_hip}{adductor muscles of the hip}, and the \href{https://en.wikipedia.org/wiki/Abdominal_external_oblique_muscle}{external}, and \href{https://en.wikipedia.org/wiki/Abdominal_internal_oblique_muscle}{internal obliques}.
    \item Secondary: gluteus maximus (glutes), quadriceps (quads), and \href{https://en.wikipedia.org/wiki/Hamstring}{hamstrings}.
\end{itemize}

\subsection{World records}
\textit{Guinness World Records} lists the record for longest duration of a front plank, resting on elbows, as 8 hours, 15 minutes and 15 seconds,[8][9] set by veteran \href{https://en.wikipedia.org/wiki/United_States_Marine_Corps}{Marine} officer George Hood on Feb 20, 2020.[10][11][12][13]

Hood also completed a record attempt in June 2018, holding a plank for 10 hours, 10 minutes and 10 seconds as well as the most cumulative plank time in a 24-hour period of 18 hours, 10 minutes and 10 seconds.[14]

%
The longest time in an elbow plank:
\begin{itemize}
    \item by a woman is 4 hours, 19 minutes and 55 seconds by Dana Glowacka (Canada) in May 2019.[15]
    \item by a woman with a 60-lb pack is 17 min and 26 sec by Eva Bulzomi (USA) in July 2013.[16]
    \item with a 100-lb pack is 17 min 02 sec by Silehm Boussehaba (France) in December 2018.[17]
    \item with a 200-lb pack is 4 min 2 sec by Silehm Boussehaba (France) in March 2018.[18]
\end{itemize}
The longest single-arm plank while balancing on medicine balls is 47.54 seconds, by William Borger (Canada) in Oct 2016.[19]

\subsection{Gallery}
\textsf{Extended plank. Side plank step by step. Modified side plank. \href{https://en.wikipedia.org/wiki/Medicine_ball}{Medicine ball} plank.}

\subsection{See also}
\begin{itemize}
    \item \href{https://en.wikipedia.org/wiki/Chaturanga_Dandasana}{Chaturanga Dandasana} - yoga low plank
    \item \href{https://en.wikipedia.org/wiki/Crunch_(exercise)}{Crunch (exercise)}
    \item \href{https://en.wikipedia.org/wiki/Planche_(exercise)}{Planche (exercise)}
    \item \href{https://en.wikipedia.org/wiki/Sit-up}{Sit-up}
    \item \href{https://en.wikipedia.org/wiki/Vasishtasana}{Vasishtasana} - yoga side plank\hfill$\square$
\end{itemize}

%------------------------------------------------------------------------------%

\section{\href{https://en.wikipedia.org/wiki/Skipping_rope}{Wikipedia\texttt{/}Skipping Rope}}
\textsf{A Ghanaian boy playing with a skipping rope.}

A \textit{skipping rope} (\href{https://en.wikipedia.org/wiki/British_English}{British English}) or \textit{jump rope} (\href{https://en.wikipedia.org/wiki/American_English}{American English}) is a tool used in the sport of skipping/jump rope where 1 or more participants jump over a \href{https://en.wikipedia.org/wiki/Rope}{rope} swung so that it passes under their feet and over their heads.

There are multiple subsets of skipping/jump rope, including single freestyle, single speed, pairs, 3-person speed (Double Dutch), and 3-person freestyle (Double Dutch freestyle).

%
There are a few major organizations that support jump rope as a sport.

Often separated by sex and age, events include hundreds of competitive teams all around the world.

In the US, schools rarely have jump rope teams, and states do not sanction official events for high school or elementary school.

In freestyle events, jumpers use a variety of basic and advanced techniques in a routine of 1 minute, which is judged by a head judge, content judges, and performance judges.

In speed events, a jumper alternates their feet with the rope going around the jumper every time 1 of their feet hits the ground for 30 seconds, 1 minute, or 3 minutes.

The jumper is judged on the number of times the right foot touches the ground in those times.

\textsf{Boy jumping a \textit{long rope} in Virginia.}

\textsf{A child playing with a skipping rope in Japan.}

\subsection{History}
\textsf{1800 illustration of a woman with a skipping rope.}

Explorers reported seeing aborigines jumping with vines in the 16th century.

European boys started jumping rope in the early 17th century.

The activity was considered indecent for girls because they might show their ankles.

Girls began to jump rope in the 18th century,[1] adding \href{https://en.wikipedia.org/wiki/Skipping-rope_rhyme}{skipping chants}, owning the rope, controlling the game, and deciding who may participate.[2]

%
In the United States, domination of the activity by girls occurred when their families moved into the cities in the late 19th century.

There, they found sidewalks and other smooth surfaces conducive to jumping rope, along with a host of contemporaries.[2]

\subsection{Techniques}
There are many techniques that can be used when skipping.

These can be used individually or combined in a series to create a routine.

\subsubsection{Solo participants}
For solo jumping, the participant jumps and swings the rope under their feet.

The timing of the swing is matched to the jump.

This allows them to jump the rope and establish a rhythm more successfully.

This can be contrasted with swinging the rope at the feet and jumping, which would mean they were matching the jump to the swing.

This makes it harder to jump the rope and establish a rhythm.

\paragraph{Basic jump or easy jump.} Jump with both feet slightly apart over the rope.

Beginners usually master this technique 1st before moving onto more advanced techniques.

\textsf{Basic jump technique.}

\paragraph{Alternate foot jump (speed step).} Use alternate feet to jump off the ground.

This technique can be used to effectively double the number of jumps per minute as compared to the above technique.

This step can be used for speed events.

\textsf{Alternate foot jump technique.}

\paragraph{Criss-cross.} Also known as crossover or cross arms.

Perform the basic jump whilst crossing arms in front of the body.

\textsf{Criss-cross technique.}

\paragraph{Side Swing.} The rope is passed by the side of the participant's body without jumping it.

\paragraph{EB (front-back cross or sailor).} Perform the criss-cross whilst crossing 1 arm behind the back.

\paragraph{Double under.} A high basic jump, turning the rope twice under the feet.

Turning the rope 3 times is called a \textit{triple under}.

In competitions, participants may attempt quadruple (quads) and quintuple under (quins) using the same method.

\paragraph{Boxer jump rope.} 1 foot is positioned slightly forward and 1 foot slightly back.

The person positions their bodyweight primarily over their front foot, with the back foot acting as a stabilizer.

From this stance, the person jumps up several times (often 2-3 times) before switching their stance, so the front foot becomes the back foot, and the back foot becomes the front foot.

And so forth.

An advantage of this technique is that it allows the back leg a brief rest.

So while both feet are still used in the jump, a person may find they can skip for longer than if they were using the basic 2-footed technique.

\paragraph{Toad.} Perform the criss-cross with 1 arm crossing under the opposite leg from the inside. 

\paragraph{Leg over/Crougar.} A basic jump with 1 arm hooked under the adjacent leg.

\textsf{Leg over technique.}

\paragraph{Awesome Annie.} Also known as Awesome Anna or swish.

Alternates between a leg over and a toad without a jump in between.

\paragraph{Inverse toad.} Perform the toad whilst 1 arm crosses the adjacent leg from the outside.

\paragraph{Elephant.} A cross between the inverse toad and the toad, with both arms crossing under one leg.

\paragraph{Frog or Donkey kick.} The participant does a \href{https://en.wikipedia.org/wiki/Handstand}{handstand}, returns to their feet, and turns the rope under them.

A more advanced version turns the rope during the return to the ground.

\paragraph{TJ.} A triple-under where the 1st `jump' is a side swing, the middle jump is a toad and the final jump in the open.

\subsubsection{Competition techniques}
\textsf{Advanced competition technique.}

In competitions, participants are required to demonstrate competence using specific techniques.

The selection depends on the judging system and the country in which the tournament is held.

\textsf{\href{https://en.wikipedia.org/wiki/Double_Dutch_(jump_rope)}{Double Dutch} - Competition during a \href{https://en.wikipedia.org/wiki/Steel_beach_picnic}{steel beach picnic} on the ship \href{https://en.wikipedia.org/wiki/USS_Saipan_(LHA-2)}{USS \textit{Saipan} (LHA-2)}.}

\subsection{Health effects}
Skipping may be used as a \href{https://en.wikipedia.org/wiki/Cardiovascular}{cardiovascular} workout, similar to \href{https://en.wikipedia.org/wiki/Jogging}{jogging} or \href{https://en.wikipedia.org/wiki/Cycling}{bicycle riding}, and has a high \href{https://en.wikipedia.org/wiki/Metabolic_equivalent}{MET} or \href{https://en.wikipedia.org/wiki/Exercise_intensity#Intensity_Levels}{intensity level}.

This \href{https://en.wikipedia.org/wiki/Aerobic_exercise}{aerobic exercise} can achieve a ``burn rate'' of up to 700 to over 1200 \href{https://en.wikipedia.org/wiki/Calorie}{calories} per hour of vigorous activity, with about 0.1 to nearly 1.1 calories consumed per jump, mainly depending upon the speed and intensity of jumps and leg foldings.

10 minutes of skipping are roughly the equivalent of running an 8-minute mile.

Skipping for 15--20 minutes is enough to burn off the calories from a candy bar and is equivalent to 45--60 minutes of running, depending upon the intensity of jumps and leg swings.

Many professional trainers, fitness experts, and professional fighters greatly recommend skipping for burning fat over any other alternative exercises like running and jogging.[3][4]

%
Weighted skipping ropes are available for such athletes to increase the difficulty and effectiveness of such exercise.

Individuals or groups can participate in the exercise, and learning proper techniques is relatively simple compared to many other athletic activities.

The exercise is also appropriate for a wide range of ages and fitness levels.

%
Skipping grew in popularity in 2020 when gyms closed or people stayed home due to \href{https://en.wikipedia.org/wiki/COVID-19_lockdowns}{Coronavirus restrictions} across the world.[5]

These workouts can be done at home and do not require a lot of equipment.

\subsection{Competition}

\subsubsection{International}
There are 2 main world organizations: International Rope Skipping Federation (FISAC-IRSF), and the World Jump Rope Federation (WJRF).

There have been 11 World Championships on every alternate year by (FISAC), with the most recent held in Shanghai, China.

There have been 7 World Jump Rope Championships held every year by (WJRF); the most recent taking place in Orlando, Florida.

Other locations of this championship include Washington DC, France, and Portugal.

%
In 2018 the International Rope Skipping Federation (FISAC-IRSF) and World Jump Rope Federation (WJRF) announced a merged organization called International Jump Rope Union.[6]

The International Jump Rope Union (IJRU) has become the 10th International Federation to gain GAISF Observer status.

The decision was taken by the Global Association of International Sports Federations (GAISF) Council, which met during SportAccord in Bangkok.

Observer status is the 1st step on a clear pathway for new International Federations towards the top of the Olympic Family pyramid.

Those who wish to proceed will be assisted by GAISF, leading them into full GAISF membership through the Alliance of Independent Recognized Members of Sport (AIMS), and the Association of IOC Recognized International Sports Federations (ARISF).

%
In 2019 the International Rope Skipping Organizations IRSO re-emerged and reactivated its activities.

The organization is headed by the founder of Rope Skipping/Jump Rope sport Richard Cendali.

IRSO[7] disagree and was not happy the way both organizations International Rope Skipping Federation and World Jump Rope Federation ignored several long-standing organizations in this merger.

Various organizations that have long-standing for the development of the sport but eft out of the merger came under IRSO under the leadership of Richard Cendali.

USA Jump Rope Federation and newly formed Asian Rope Skipping Association also joined IRSO and decided to host their World Championship in conjunction with AAU.

\subsubsection{World Inter School}
The 1st World Inter-School Rope Skipping Championship [8] was held at Dubai, Nov 2015.

The 2nd World Inter-School Rope Skipping Championship was held at Eger, Hungary.

The Championship was organized by World Inter School Rope Skipping Organisation (WIRSO).[9]

2nd, 3rd and 4th [10] World Inter-School championships held in Hungary 2017, Hong Kong 2018 and Belgium 2019 respectively. 

\subsubsection{United States}
Historically in the United States there were 2 competing jump rope organizations: the International Rope Skipping Organization (IRSO), and the World Rope Skipping Federation (WRSF).

IRSO focused on \href{https://en.wikipedia.org/wiki/Stunt}{stunt}-oriented and gymnastic/athletic type moves, while the WRSF appreciated the \href{https://en.wikipedia.org/wiki/Aesthetics}{aesthetics} and form of the exercise.

In 1995 these 2 organizations merged to form the United States Amateur Jump Rope Federation (now, \href{http://www.usajumprope.org/}{USA Jump Rope}).

\href{http://www.usajumprope.org/}{USA Jump Rope} hosts annual national tournaments, as well as camps, workshops, and clinics on instruction.

Jump rope is also part of the \href{https://en.wikipedia.org/wiki/Amateur_Athletic_Union}{Amateur Athletic Union} and participates in their annual \href{https://en.wikipedia.org/wiki/AAU_Junior_Olympic_Games}{AAU Junior Olympic Games}.[11]

%
Speed jump ropes are made from a thin vinyl cord.

They are best for indoor use, because they will wear down fast on concrete or other harsh surfaces.

The beaded ropes make rhythmic jumping very easy, because the jumper can hear the beads hitting the ground and strive for a rhythmic pattern.

The leather jump rope tangles less than the speed rope.

\subsection{See also}
\begin{itemize}
    \item \href{https://en.wikipedia.org/wiki/Chinese_jump_rope}{Chinese jump rope}\hfill$\square$
\end{itemize}

%------------------------------------------------------------------------------%

\printbibliography[heading=bibintoc]
\end{document}